\section*{Erd\H{o}s Problem 908: the measurable conclusion and a wording correction}

\paragraph{Literal statement and the primary theorem.}
The current problem page assumes that every difference
\(\Delta_t f\) is measurable, but asks for
\(f=g+H+r\) with \(g\) \emph{continuous}, \(H\) additive,
and \(\Delta_t r=0\) almost everywhere for each fixed \(t\).
That literal continuous-\(g\) statement is false.
Laczkovich's original Theorem 3 instead has \(g\)
\emph{Lebesgue measurable}. We first give the counterexample
to the literal wording, then reconstruct the corrected
theorem relative to two explicitly stated preparatory
results from the same paper.

\paragraph{A counterexample to continuity of \(g\).}
Take \(f(x)=\mathbf1_{[0,\infty)}(x)\). It is measurable,
so all its differences are measurable.
Its difference with shift one equals one on \((-1,0)\)
and zero on \((-\infty,-1)\cup(0,\infty)\).
If the displayed continuous-\(g\) decomposition existed,
then for almost every \(x\)
\[
 \Delta_1 f(x)=g(x+1)-g(x)+H(1).
\]
The right side is continuous. A continuous function
equal almost everywhere to a constant on an open
interval equals that constant everywhere there:
otherwise a neighbourhood of positive measure
would contradict the almost-everywhere equality.
It would therefore equal zero just to the left of
\(-1\) and one just to its right, an impossibility.
This argument does not assume that \(r\) is measurable
and uses only the null set associated with shift one.

\paragraph{Two explicit inputs for the corrected theorem.}
On measurable 1-periodic functions modulo equality
almost everywhere, use the translation-invariant metric
\[
 \|u\|_0=\int_0^1\min(1,|u(x)|)\,dx,\qquad
 s(u)=\inf_{a\in\mathbb R}\|u-a\|_0.
\]
Convergence for this metric is convergence in measure.
We retain the following two results of Laczkovich as
external inputs, stated in their equivalent
convergence-in-measure form.
\begin{enumerate}
\item Theorem 2: if \(u_j\) are measurable 1-periodic
functions and \(\Delta_tu_j\to0\) in measure for every
fixed real \(t\), then \(s(u_j)\to0\).
\item Theorem 1: if \(c:\mathbb R\to\mathbb R\),
\(c(0)=0\), and
\[
 c(h+k)-c(h)-c(k)\longrightarrow0\quad((h,k)\to(0,0)),
\]
then \(c=H+u\), where \(H\) is additive and \(u(h)\to0\)
as \(h\to0\), with \(u(0)=0\).
\end{enumerate}
Their proofs are not reproduced here. Standard
translation continuity in measure, Fubini's theorem,
and the Borel--Cantelli lemma will also be used.

\paragraph{Reduction to a periodic function.}
All real shifts have measurable differences, since
\(\Delta_{-t}f(x)=-\Delta_tf(x-t)\).
Define \(F\) by \(F(x)=f(x)\) on \([0,1)\) and
extend it with period one.
On each interval \([j,j+1)\),
\(f(x)-F(x)=f(x)-f(x-j)\) is measurable.
Thus \(f-F\) is measurable, and every difference of
\(F\) is measurable. It suffices to decompose \(F\).
Put \(D_h(x)=F(x+h)-F(x)\).

For \(h_j\to0\) and fixed \(t\),
\[
 \Delta_tD_{h_j}(x)=\Delta_{h_j}D_t(x)\longrightarrow0
 \quad\hbox{in measure}.
\]
The first external input gives \(s(D_h)\to0\) as \(h\to0\).
Choose a minimizing constant \(c(h)\), so that
\(\|D_h-c(h)\|_0=s(D_h)\), choosing periodically
and with \(c(0)=0\).
Such a minimum exists: the metric is continuous in the
constant and tends to one as its absolute value tends
to infinity; if the infimum is one, every constant minimizes.

The identity
\[
 D_{h+k}(x)=D_h(x+k)+D_k(x)
\]
and the triangle inequality imply
\[
 \min(1,|c(h+k)-c(h)-c(k)|)
 \leq s(D_{h+k})+s(D_h)+s(D_k)\longrightarrow0.
\]
Apply the second external input to obtain \(c=H+u\).
Replacing \(H(x)\) by \(H(x)-xH(1)\) and \(u(x)\)
by \(u(x)+xH(1)\), we may suppose \(H(1)=0\).

\paragraph{A jointly measurable representative.}
Set
\[
 K(x,h)=F(x+h)-F(x)-H(h).
\]
For fixed \(h\), this is measurable and 1-periodic in \(x\).
Moreover
\[
 \|K(\,\cdot\,,h)\|_0
 \leq s(D_h)+\min(1,|u(h)|)\longrightarrow0
 \quad(h\to0).
\]
Its cocycle identity gives
\[
 \|K(\,\cdot\,,h)-K(\,\cdot\,,h')\|_0
 =\|K(\,\cdot\,,h-h')\|_0.
\]
Choose positive meshes \(\eta_j\to0\) so small that
the last norm is below \(2^{-3j}\) whenever
\(|h-h'|\leq\eta_j\), and define
\[
 G_j(x,h)=K(x,\eta_j\lfloor h/\eta_j\rfloor).
\]
Each \(G_j\) is jointly Lebesgue measurable, and its
section as a function of \(h\) is Borel for every \(x\).
For fixed \(h\),
\[
 \lambda\{x\in[0,1]:|G_j(x,h)-K(x,h)|\geq2^{-j}\}
 \leq2^{-2j}.
\]
Borel--Cantelli shows that \(G_j(x,h)\to K(x,h)\)
for almost every \(x\), for each fixed \(h\).
Define \(G(x,h)\) as the finite limit when it exists,
and zero otherwise. Then \(G\) is jointly measurable,
each \(G(x,\cdot)\) is Borel, and
\[
 R(x,h):=K(x,h)-G(x,h)=0
 \quad\hbox{for a.e. }x,\quad\hbox{for every fixed }h.
\tag{1}
\]
Periodicity in \(x\) extends this assertion from \([0,1]\)
to the whole line.

\paragraph{Extracting the almost-invariant remainder.}
Although joint measurability of \(R\) is not presumed,
its defect is jointly measurable:
\[
 \begin{split}
 L(x,y,z)
 &:=R(x,y+z)-R(x+y,z)-R(x,y)\\
 &=-G(x,y+z)+G(x+y,z)+G(x,y).
 \end{split}
\]
For fixed \(y,z\), (1) makes \(L(x,y,z)=0\) for a.e. \(x\).
Fubini therefore supplies \(x_0\) such that
\(L(x_0,y,z)=0\) for almost every pair \((y,z)\).
For almost every \(z\), this and (1) give
\[
 R(x_0,y+z)-R(x_0,y)=0\quad\hbox{for a.e. }y.
\tag{2}
\]
Let \(Z\) be the full-measure set of such \(z\).
For any real \(t\), choose \(z_1,z_2\in Z\) with
\(t=z_1+z_2\); this is possible because
\(Z\cap(t-Z)\) has full measure.
Adding the two translated instances of (2) proves
the same equality for shift \(t\).
The exceptional null set may depend on \(t\), exactly
as required in the question.

Now put
\[
 r(x)=R(x_0,x-x_0),\qquad
 g_0(x)=G(x_0,x-x_0)+F(x_0)-H(x_0).
\]
By construction \(g_0\) is measurable,
\(F=g_0+H+r\), and \(\Delta_t r=0\) a.e. for every \(t\).
Finally \(g=g_0+(f-F)\) is measurable and
\(f=g+H+r\), proving the corrected theorem.

\paragraph{Sources and assessment.}
The exact primary source is M. Laczkovich,
\emph{Functions with measurable differences},
Acta Math. Acad. Sci. Hungar. 35 (1980), 217--235,
especially Theorems 1--3 on pages 218--227:
\url{https://real-j.mtak.hu/7443/1/MTA_ActaMathHung_35.pdf}.
The erroneous continuous-\(g\) wording also appears
in Erd\H{o}s's 1982 exposition, page 76,
\url{https://users.renyi.hu/~p_erdos/1982-33.pdf};
it is not merely being attributed to a website transcription.
The current page
\url{https://www.erdosproblems.com/908} was checked
24 September 2026.
The literal counterexample is elementary and complete.
The corrected affirmative theorem is completely deduced
from the two declared preparatory inputs, following
Laczkovich's proof. No novelty or local formal
verification is claimed.

\paragraph{Completion estimate.}
\textbf{COMPLETION: 100\%}
